\section{Retained leading statements and their proof routes}
\label{app:v49-retained-leading}
The following statements reproduce the preceding leading theorem sequence.
Their detailed proofs remain in the unchanged core sections. All references
use labels rather than the historical order of the front matter.

\begin{leadtheorem}[The stationary microscopic local law]
\label{thm:intro-stationary-microscopic}
There is a continuous uniformly positive definite matrix $\mathcal V_R$
such that, for every $M<\infty$,
\begin{equation}\label{eq:intro-stationary-local-law}
 \sup_{\substack{R\in I,\ k\in\Z^2,\ m\in\mathbb N\\
                     |\xi_{t,R}(k,m)|\le M}}
 \left|t^{3/2}\mathbb P_R\{W_R(t)=k,C_R(t)=m\}
              -g_{\mathcal V_R}(\xi_{t,R}(k,m))\right|\longrightarrow0,
\end{equation}
where $I=[0.45,0.47]$ and
$\xi_{t,R}(k,m)=(k/\sqrt t,(m-t/\bar\tau_R)/\sqrt t)$.
The covariance is that of the stationary physical functional limit.
For bounded regular initial and terminal flight-state functions
$\chi_R,\psi_R$, the same formula holds with the probability replaced
by its $\chi_R\psi_R$-weighted expectation and the Gaussian multiplied
by $\mu_R(\chi_R)\mu_R(\psi_R)$.

For every fixed $P>0$, the entire signed complement in
\eqref{eq:stationary-middle-complement}, with $B_m=m^{P+3/2}$ and
physical central radius $2m^{-99/200}$, satisfies
$t^{3/2}\mathfrak M_{m,R,B_m}(k,t)=o(1)$ on these central sets.
The positive likelihood for the same microscopic event then approximates
its normalized conditional measure on the original trajectory space
in total variation with error $O(m^{-P})$. This comparison is uniform
against all bounded measurable trajectory tests.
\end{leadtheorem}
\begin{proof}
Theorems~\ref{thm:stationary-microscopic-LLT} and
\ref{thm:stationary-endpoint-selected-LLT}, followed by
Corollaries~\ref{cor:stationary-middle-vanishes} and
\ref{cor:normalized-physical-band}, prove the assertions. The regular
endpoint class is specified before
Theorem~\ref{thm:stationary-endpoint-selected-LLT}.
\end{proof}

\begin{leadtheorem}[An action principle beyond circular tables]
\label{thm:intro-action-family}
Place one oriented ellipse at every point of the same triangular lattice,
with semiaxes $a,b\in[0.45,0.47]$ and orientation $\theta\in\R/2\pi\Z$.
The stationary local law \eqref{eq:intro-stationary-local-law}, its regular
endpoint-selected form and its normalized positive-likelihood comparison
hold uniformly in $(a,b,\theta)$. The mean roof is
\[
 \bar\tau_{a,b,\theta}=\frac{\pi(\sqrt3/2-\pi ab)}{L_{a,b}},
\]
where $L_{a,b}$ is the perimeter. The limiting physical covariance is
continuous and uniformly positive definite, including at $a=b$.

More generally, these conclusions hold for compact families satisfying
the separated geometry, local-space, action, finite-cover mixing,
parameter continuity, preceding-flight and regular observation conditions in Section~\ref{sec:compact-family-local-principle}.
For these families the physical phase equation
$q\circ T=e^{i(u\cdot\kappa+s+b_0\tau)}q$ has a measurable circle-valued
solution only when $b_0=0$, $u\in(2\pi\Z)^2$ and $s\in2\pi\Z$.
Neither a collision rotation nor a quantitative bound on growing-band
spectral constants is assumed.
\end{leadtheorem}
\begin{proof}
Theorem~\ref{thm:v35-rotation-free-rigidity},
Theorem~\ref{thm:v35-action-family-LLT},
Lemma~\ref{lem:v35-ellipse-geometry} and
Corollary~\ref{cor:v35-family-posterior} give the assertions.
The original circular theorem is the restriction $a=b=R$; the signed
middle identity in its statement retains its previously specified
circular record and cutoff.
\end{proof}


\begin{leadtheorem}[Local constraints on the actual return record]
\label{thm:intro-v37-return-window}
For the original circular family let $c=\nu(Y_R^*)=91/(10000\pi)$.
There are continuous $\beta_R\in\R^{1\times3}$ and
$\sigma_{A,R}^2>0$, with uniform positive variance bounds, such that,
for every fixed bounded interval $J$ of positive length and fixed
$a<b$, uniformly for $R\in I$ and $|z|\le M$,
\begin{align}\label{eq:intro-v37-return-window}
 &m^{3/2}\sum_{\substack{1\le n\le m\\a\le(n-mc)/\sqrt m\le b}}
 \nu_R^*\{K_{n,R}=k,\ N_{n,R}=m,\ T_{n,R}-t\in J\}\notag\\
 &\hspace{12mm}
 -c|J|g_{\Sigma_R}(z)
       \int_a^b g_{\sigma_{A,R}^2}(y-\beta_Rz)\,dy
                           \longrightarrow0,\qquad
 z=\left(\frac{k}{\sqrt m},\frac{t-m\bar\tau_R}{\sqrt m}\right).
\end{align}
Here $\Sigma_R$ is the collision displacement--roof covariance.
The Schur data $\beta_R,\sigma_{A,R}^2$ are those of the joint
covariance after adjoining the true section occupation. The summands
are probabilities of the original return records; neither the section
nor the event is replaced. The probability sum is bounded below by
$c_{M,a,b,J}m^{-3/2}$ for all sufficiently large $m$.
\end{leadtheorem}
\begin{proof}
Lemma~\ref{lem:v37-section-multiplier} realizes the exact occupation.
Theorem~\ref{thm:v37-occupation-local-central} proves its mixed
local--central law. Proposition~\ref{prop:v37-exact-return-disintegration}
and Theorem~\ref{thm:v37-return-window} give the statement; equation
\eqref{eq:v37-return-gaussian-agreement} identifies its amplitude with
the Gaussian normalization of the original induced covariance $D_R$.
\end{proof}


\begin{leadtheorem}[Fine resolution of the actual return index]
\label{thm:intro-v38-fine-returns}
For the original circular family and every bounded interval $J$,
\[
 \nu_R^*\{K_{n,R}=k,N_{n,R}=m,T_{n,R}-t\in J\}
       \le C(1+|J|)m^{-2}\qquad(m\ge n\ge1),
\]
uniformly in the radius and all targets. Let $J$ now have fixed positive
length. Let $\mathcal B_m$ consist of $H_m$ consecutive actual return
indices, with center $\ell_m$, where $H_m\to\infty$ and
$H_m/\sqrt m\to0$. Put
\[
 Z_{m,R}=\left(\frac{k}{\sqrt m},
       \frac{t-m\bar\tau_R}{\sqrt m},
       \frac{\ell_m-mc}{\sqrt m}\right),\qquad
 \Omega_R=cL_RD_RL_R^{\mathsf T},
\]
with $L_R$ specified in \eqref{eq:v37-covariance-change}. Then
\begin{align}\label{eq:intro-v38-fine-returns}
 \frac{m^2}{H_m}\sum_{n\in\mathcal B_m}
 \nu_R^*\{K_{n,R}=k,N_{n,R}=m,T_{n,R}-t\in J\}
       =c|J|g_{\Omega_R}(Z_{m,R})+o(1),
\end{align}
uniformly on central compact sets and over widths between any
$h_m\to\infty$ and $u_m\sqrt m$ with $u_m\to0$.
The resulting denominator is at least $c_{M,J}H_mm^{-2}$.
No lower polynomial growth restriction is imposed on $H_m$.
The upper bound applies to one index; the Gaussian asymptotic in this
statement uses a diverging block, and the roof interval is not replaced by a pointwise
density. The convergence is qualitative, with no asserted polynomial
rate for the $o(1)$ term.
\end{leadtheorem}
\begin{proof}
Theorem~\ref{thm:v38-fine-return-windows} gives the assertions.
Corollary~\ref{cor:v38-return-normalization} writes the same limit in
the original $D_R$ normalization. The explicit two-band budget is
\eqref{eq:v38-two-band-budget};
Proposition~\ref{prop:v38-exact-index-remainder} isolates the remaining
single-index occupation-torus inverse.
\end{proof}

\begin{leadtheorem}[An exact return index with its arithmetic coefficient]
\label{thm:intro-arithmetic-return}
For the original circular table and the unchanged section $Y_R^*$,
all measurable occupation resonances have zero roof coordinate. Their
group is finite cyclic, of order $d_R\le D$ uniformly in $R$. At each
fixed $R$, let $a_R^{\rm ar}\in(\Z/d_R\Z)^2$, $h_R^{\rm ar}\in\Z/d_R\Z$ and
$w_{R,l}\ge0$ be the residue data in
\eqref{eq:v40-section-phase-masses}, with $\sum_lw_{R,l}=c$.
Put
\[
 \mathfrak a_R(k,n,m)=\frac{d_R}{c^2}
   \sum_{l\bmod d_R}w_{R,l}w_{R,l+a_R^{\rm ar}\cdot k+n+h_R^{\rm ar}m}.
\]
For every fixed bounded interval $J$ of positive length and each
fixed radius,
\begin{equation}\label{eq:intro-arithmetic-return}
 n^2\nu_R^*\{K_{n,R}=k,N_{n,R}=m,T_{n,R}-t\in J\}
   =|J|\mathfrak a_R(k,n,m)g_{D_R}(V_n)+o_R(1),
\end{equation}
where $V_n=n^{-1/2}(k_1,k_2,m-n/c,t-n\bar\tau_R/c)$ and the error
is uniform over its compact sets. The coefficient is nonnegative,
its average over $d_R$ consecutive return indices is one, and a
zero coefficient makes the exact event null. There is a positive
microscopic denominator on each positive residue class.

For every fixed roof band, the complete fixed-count inverse also has
a finite spectral reduction uniform over all $R\in I$, as stated in
\eqref{eq:v40-fixed-spectral-inverse}. This uniform formula retains
near-resonant branches through parameter changes. No unmodulated
radius-uniform singleton law or pointwise roof-density statement is
included in \eqref{eq:intro-arithmetic-return}.
\end{leadtheorem}
\begin{proof}
Apply Theorems~\ref{thm:v40-uniform-resonant-reduction},
\ref{thm:v40-no-roof-resonance} and
\ref{thm:v40-arithmetic-return-LLT}.
\end{proof}

\begin{leadtheorem}[Uniform transitions and the exact-event return bridge]
\label{thm:intro-v41-bridge}
Let $J$ be a fixed bounded interval of positive length and put
\[
 E=\{K_{n,R}=k,N_{n,R}=m,T_{n,R}-t\in J\},\qquad
 Z=m^{-1/2}(k_1,k_2,t-m\bar\tau_R,n-mc).
\]
The probability of this exact return event satisfies
\[
 m^2\nu_R^*(E)=|J|\mathcal L_{m,R}(k_1,k_2,t,n)+o(1)
\]
uniformly in $R$ and the targets, where $\mathcal L_{m,R}$ is the
finite Gaussian expression in \eqref{eq:v41-transition-kernel}.
It retains the damping and drift of branches approaching the unit
circle; at fixed $R$ it reduces to the arithmetic Gaussian in Theorem~\ref{thm:intro-arithmetic-return}.

Let $\mathcal Y_{n,R}$ be the polygonal path with values
$n^{-1/2}(J_{l,R}-(l/n)J_{n,R})$ at $l/n$.
For every $M,\epsilon>0$, uniformly on $|Z|\le M$ and
$m^2\nu_R^*(E)\ge\epsilon$,
\[
 d_{\rm BL}\left(\operatorname{Law}(\mathcal Y_{n,R}\mid E),
                         \operatorname{Law}(\mathbb B_{D_R})\right)
                                      \longrightarrow0.
\]
Here $\mathbb B_{D_R}$ is the centered continuous Gaussian bridge with
covariance $(s\wedge t-st)D_R$. At each fixed radius the mass condition
holds on every positive arithmetic residue class. There is also a fixed
finite integer $L$, independent of $m$ and $R$, such that
\[
 \frac{m^2}{L}\sum_{j=0}^{L-1}
 \nu_R^*\{K_{n+j,R}=k,N_{n+j,R}=m,T_{n+j,R}-t\in J\}
                         =c|J|g_{\Omega_R}(Z)+o(1)
\]
uniformly on central compact sets. This last packet assertion does not
replace the single-index event used in the bridge statement.
\end{leadtheorem}
\begin{proof}
Apply Theorems~\ref{thm:v41-uniform-transition},
\ref{thm:v41-actual-return-bridge} and \ref{thm:v41-finite-packet}.
Theorem~\ref{thm:v41-zero-residue-criterion} characterizes the additional
condition for the radius-uniform unmodulated single-index interval law.
\end{proof}

\paragraph{Arithmetic and cancellation at an exact index.}
The new theorem estimates the full occupation torus at a prescribed
collision count. Section~\ref{sec:v40-occupation-powers} refines each
regular inverse curve by at most a linear number of occupation cuts.
The resulting polynomial losses are absorbed in the three collision
norm inequalities. Section~\ref{sec:v40-holonomy-residues} proves
that every occupation holonomy differs from its roof holonomy by an
integer; countability then removes nonzero roof resonances. It also
calculates the true section-to-section coefficient at every remaining
scalar eigenvalue. Section~\ref{sec:v40-fixed-count-local} performs
the fixed-count inverse and sums those residues by finite Fourier
orthogonality.

The earlier source spectral and Abel identities in
Section~\ref{sec:v39-section-source} remain exact. They are not used
to infer decay of a fixed coefficient. The signed complement from
\eqref{eq:v38-fixed-torus-remainder} is now evaluated at fixed $R$:
its main term is proportional to $\mathfrak a_R-1$, not automatically
zero. A fixed interval is not a pointwise roof density. Section~\ref{sec:v41-transition} now evaluates arithmetic transitions
uniformly by retaining their damped spectral branches. Sections~
\ref{sec:v41-bridges} and~\ref{sec:v41-residue-filters} give the
exact-event path bridge and a finite-packet consequence. Triviality of
the section residues and the common pointwise raw correction remain
separate requirements for the original full raw theorem.

\begin{leadtheorem}[Regular critical packets and arithmetic endpoint laws]
\label{thm:intro-v42-critical-packet}
For the original circular family and the actual section, every finite
collision cutoff $L\ge n$ admits the exact mixed-density decomposition
\[
 p_{n,R}=e_{n,R}^{L,\varepsilon}+q_{n,R}^{L,\varepsilon}
                                +r_{n,R}^{L,\varepsilon},\qquad r\ge0.
\]
The finite edge sum $e$ retains every guarded interior critical word;
no initial normal strip is discarded. Each edge has the explicit form
$\one_{s\ge0}e^{-s}\{a_0+(a_0+a_1)s\}$, where $a_0,a_1$ are the
intrinsic right roof jets. Its coefficients stabilize on each word with positive guard margins
as the guards are removed. Uniformly in $R$,
\begin{align*}
 \|r\|_{L^1(\Z^3\times\R)}
   &\le C\{e^{an-c_0L}+(L+1)\varepsilon^{1/16}\},\\
 \sum_z(|a_{0,z}|+|a_{1,z}|)+\sum_\ell\|q(\ell,\cdot)\|_{W^{2,1}}
   &\le \exp\{C(L+1)\log(C/\varepsilon)\}.
\end{align*}
The residual inverse truncated at roof frequency $B$ has pointwise
error at most the last bound divided by $\pi B$. The uncut source has
in addition the nonnegative pointwise error $r(\ell,t)$; its small
mass alone is not a uniform pointwise estimate.

At each fixed radius, let $w_l$ be the actual section phase masses and
$r_0=r_R(k,n,m)$ its arithmetic residue. Under the unchanged central
exact-index event with a fixed positive roof interval and
$\sum_lw_lw_{l+r_0}>0$, the endpoint phase pair has limiting masses
\[
 \pi_{R,r_0}(l,l')=
 \one_{\{l'=l+r_0\}}\frac{w_lw_{l+r_0}}{\sum_hw_hw_{h+r_0}}.
\]
These phases are asymptotically independent of the actual pinned return
Gaussian bridge. No regularity of the measurable phase indicators as
strong multipliers is assumed.
\end{leadtheorem}
\begin{proof}
Theorem~\ref{thm:v42-regular-critical-packet},
Corollary~\ref{cor:v42-pointwise-packet},
Theorem~\ref{thm:v42-endpoint-posterior} and
Corollary~\ref{cor:v42-phase-bridge-product} give the assertions.
\end{proof}

\begin{leadtheorem}[Deconcentration of regular critical edges]
\label{thm:intro-v45-critical-cluster}
For $0<\varepsilon<1/10$ let $\mathcal J^\varepsilon_{n,k,m,R}$ be the
positive measure of intrinsic critical jump coefficients in
\eqref{eq:v45-critical-coefficient-measure}, using the exponentially
relaxed margins of Section~\ref{sec:v45-critical-collars}. There are
$C,h_*>0$, independent of $\varepsilon$, such that for every fixed
$0<h<h_*\varepsilon^6$,
\[
 \limsup_{m\to\infty}\sup_{R,n,k,t}
       m^2\mathcal J^\varepsilon_{n,k,m,R}([t-h,t+h])\le Ch.
\]
The sum at any single critical value is therefore $o(m^{-2})$,
uniformly in the exact labels and radius for each fixed $\varepsilon$.
For every sequence $r_m\downarrow0$, the original source restricted
to the physical critical collars of roof excess less than $r_m$ has
density $p^{\varepsilon,\mathrm{crit}}$ satisfying
\[
 \sup_{R,n,k}m^2\|p^{\varepsilon,\mathrm{crit}}_{n,k,m,R}\|_\infty
       \longrightarrow0.
\]
The same conclusion holds for arbitrary bounded measurable insertions
on this restricted source. Its contribution to $p-K_B*p$ is pointwise
$o(m^{-2})$, uniformly for all $B>0$. No assertion about the full
remaining signed correction or about a prescribed shrinking margin
$\varepsilon=\varepsilon_m$ is included.
\end{leadtheorem}
\begin{proof}
Apply Theorems~\ref{thm:v45-critical-cluster} and
\ref{thm:v45-critical-source-smallness}, followed by
Corollary~\ref{cor:v45-critical-unsmoothing}.
\end{proof}

\paragraph{Arithmetic and the pointwise return target.}
For the exact return index we retain the factor $\mathfrak a_R$ at
fixed radius and the transition kernel $\mathcal L_{m,R}$ uniformly
in the radius. The unmodulated specialization requires the zero-residue
criterion; it is not imposed by notation. Theorem~\ref{thm:LLT}
remains the inherited raw inversion criterion, not a statement that
its remaining pointwise hypotheses have been verified. Theorem~
\ref{thm:intro-v42-critical-packet} advances that same raw problem by
retaining, rather than removing, the regular critical source. It does
not replace the pointwise task by an existential shrinking interval.


\begin{leadtheorem}[Protected raw inversion and margin localization]
\label{thm:intro-v46-protected-inverse}
Let $\Xi_{m,R}^{\varepsilon}$ be the graded source guard in
\eqref{eq:v46-source-guard}, equal to one when all the graded margins
exceed twice their thresholds. For the original exact labels put
\[
 f^\varepsilon(t)\,dt
  =(T_{n,R})_*(\one_{\{K_{n,R}=k,\,N_{n,R}=m\}}
                   \Xi_{m,R}^{\varepsilon}\,\nu_R^*),\qquad
 b^\varepsilon=p_{n,R}(k,m,\cdot)-f^\varepsilon\ge0.
\]
For every fixed $\varepsilon>0$ and $B\ge C_2\varepsilon^{-12}$,
\begin{equation}\label{eq:intro-v46-band-modulus}
 \limsup_{m\to\infty}\sup_{R,n,k}
     m^2\|f^\varepsilon-K_B*f^\varepsilon\|_\infty
                                 \le C B^{-1/2}.
\end{equation}
The same correction tends to zero for every $B_m\to\infty$ with
$\varepsilon$ fixed, without an asserted rate in $m$. Regular endpoint
insertions satisfy the corresponding estimate with their stated
supremum and endpoint-gradient budgets. The positive remainder obeys
\[
 \sup_{m,R}\sum_{n=1}^m\sum_k\|b^\varepsilon\|_1
                                    \le C\varepsilon.
\]
With $\varepsilon(B)=AB^{-1/12}$ for a fixed sufficiently large $A$,
the full raw arithmetic local law with main term $\mathcal L_{m,R}$
is equivalent to vanishing of the normalized signed correction
$b^{\varepsilon(B)}-K_B*b^{\varepsilon(B)}$ on each central compact
set, first in collision limsup and then as $B\to\infty$.
The mass bound alone does not supply this last estimate.
\end{leadtheorem}
\begin{proof}
Apply Theorems~\ref{thm:v46-noncritical-density},
\ref{thm:v46-protected-correction} and
\ref{thm:v46-boundary-reduction}, and
Proposition~\ref{prop:v46-protection-mass}.
\end{proof}


\begin{leadtheorem}[Pointwise removal of endpoint decision boundaries]
\label{thm:intro-v47-boundary-decisions}
Fix an endpoint depth $J\ge0$ and a margin $\varepsilon>0$. Let
$G^{\varepsilon,J}_{m,R}$ be the original graded protection with the
section factors at $0,\ldots,J,m-J,\ldots,m$ omitted. Keep every
incidence and clearance factor and every other section factor. Let
$E^{\varepsilon,J}_{m,R}$ be the product of the omitted factors, and
for the original exact labels define
\[
 d^a(t)\,dt=(T_{n,R})_*
 \bigl(\one_{\{K_{n,R}=k,N_{n,R}=m\}}
       aG^{\varepsilon,J}_{m,R}(1-E^{\varepsilon,J}_{m,R})\nu_R^*\bigr).
\]
Uniformly for all measurable insertions with $|a|\le M$,
\[
 \limsup_{m\to\infty}\sup_{R,n,k,a}m^2\|d^a\|_\infty
       \le CM(J+1)\varepsilon.
\]
The same bound holds for $d^a-K_B*d^a$, uniformly over every $B>0$.
All original integer labels and the roof observable on the left are
unchanged. The finite occupation enlargement occurs only in a positive
upper comparison, not as an average replacing the target coefficient.

In particular, with $\varepsilon(B)=AB^{-1/12}$ and
$J(B)=\lfloor B^{1/24}\rfloor$, the full unweighted correction differs
from the signed correction of $1-G^{\varepsilon(B),J(B)}$ by at most
$CB^{-1/24}$ at the normalized collision-limsup scale. This residual
splits by its first bad margin into incidence, clearance and interior
section-decision sources. Its pointwise correction remains to be
estimated, with the arithmetic main term $\mathcal L_{m,R}$ retained.
\end{leadtheorem}
\begin{proof}
Apply Theorem~\ref{thm:v47-endpoint-boundary-smallness},
Proposition~\ref{prop:v47-first-defect-masses} and
Theorem~\ref{thm:v47-ordered-boundary-removal}.
\end{proof}

\begin{leadtheorem}[Every section-decision depth]
\label{thm:intro-v48-all-decisions}
Let $H_{m,R}^{\eta}$ contain all graded incidence and clearance factors,
and let $D_{m,R}^{\varepsilon}$ contain all graded section-decision
factors, as in \eqref{eq:v48-physical-decision-guards}. For the original
exact event $\Pi_{n,k,m,R}$ and any bounded measurable source insertion
$|w|\le M$, define
\[
 d^{\eta,\varepsilon,w}_{n,k,m,R}(t)\,dt
 =(T_{n,R})_*\bigl(\one_{\Pi_{n,k,m,R}}wH_{m,R}^{\eta}
                        (1-D_{m,R}^{\varepsilon})\nu_R^*\bigr).
\]
For fixed positive $\eta,\varepsilon$ sufficiently small,
\[
 \limsup_{m\to\infty}\sup_{R,n,k,\,|w|\le M}
             m^2\|d^{\eta,\varepsilon,w}_{n,k,m,R}\|_\infty
                           \le CM\varepsilon^{1/16}.
\]
The same bound holds for $d-K_B*d$, uniformly over every $B>0$.
For the positive depth telescope of \eqref{eq:v48-depth-telescope},
its tail beyond depth $J\ge-1$ satisfies the finite-count bound
\begin{align*}
 m^2\|d_{>J}^{\eta,\varepsilon,w}\|_\infty
 &\le CM\varepsilon^{1/16}(J+2)(47/53)^{(J+1)/64}\\
 &\quad+C_{\eta,\varepsilon}M m^4\rho_{\eta,\varepsilon}^{m/2},
                      \qquad \rho_{\eta,\varepsilon}<1.
\end{align*}
Consequently an increasing decision depth cannot escape this estimate
at fixed physical margins. With $\varepsilon(B)=AB^{-1/12}$ fixed
before the collision limit, the complete signed correction differs
from the incidence--clearance correction by $O(B^{-1/192})$ in
normalized essential supremum. The exact labels and arithmetic
transition main term remain unchanged.
\end{leadtheorem}
\begin{proof}
Apply Theorem~\ref{thm:v48-all-decision-smallness},
Corollary~\ref{cor:v48-decision-all-band}, and
Theorem~\ref{thm:v48-two-source-reduction}.
\end{proof}

\paragraph{Mechanism of the microscopic estimate.}
The exact age-overlap formula
\eqref{eq:exact-overlap-probability} retains the initial and terminal
cell offsets. Its two interval endpoints give a pointwise far-frequency
error $O(B^{-1})$, but that estimate alone says nothing about intermediate
frequencies. Theorem~\ref{thm:compact-collision-spectrum} supplies the
missing bounded-band dynamical mechanism for the unsmoothed collision
triple $(\kappa_1,\kappa_2,\tau-\bar\tau)$. It realizes the twist on
the image side of the collision operator. The exact action formula
controls the accumulated phase on stable curves and on the matched
curves needed for the unstable norm; physical measurable phase rigidity
excludes every nonzero peripheral frequency.

Upper and lower Fourier envelopes, applied to positive source measures,
give a joint collision local limit with both endpoint states. The order of limits is a fixed
frequency band, the collision count tending to infinity, and finally
the band tending to infinity. The original physical overlap is a
legitimate test of the resulting local measure. This evaluates the
complete probability before the Gaussian-plus-complement identity is
used to deduce the signed middle estimate. It requires neither a bound
on growing-band operator constants nor an absolute estimate for the
modulus of that middle transform.

\paragraph{Actual returns and conditioning.}
The physical local theorem does not replace the four-coordinate raw
return problem by a different record. All the return geometry,
nondegeneracy, Gaussian, marked, path and raw-inversion results remain
in the body. Their complementary formulations are collected in
Appendix~\ref{app:retained-return-statements}, with the letters A--X
used in the cross-references. The current theorem provides a direct
route to microscopic physical conditioning, using the original event
and a proved denominator; it needs no return-to-physical event
replacement. Its arbitrary-test total-variation comparison is not an
assertion that every selected path event has a positive Gaussian
amplitude. The Gaussian amplitude is evaluated here for the stated
regular endpoint class. The general weighted return-density theorem
remains separate from the exact-index, fixed-interval conditional
return bridge proved in Theorem~\ref{thm:intro-v41-bridge}. A pointwise
roof-conditioned path law is not asserted.

\paragraph{Relation to previous local limit methods.}
The theorem-level comparison in Section~\ref{sec:v35-literature-routes} distinguishes existing fixed-table local limits from the compact-family action and source estimates here. The collision theory of \cite{DZ11,DZ} and the image-side displacement
multiplier of \cite{DPZ} provide the anisotropic norm framework.
Our compact-frequency argument includes the physical roof by verifying
its action bounds in those norms; a spatial local limit theorem is not
cited as a theorem for the three-coordinate record. The endpoint form
of a mixing local limit is also natural in the flow theory of
\cite{DN}. Here uniformity in the moving family, the actual stationary
age, and the two cell offsets are treated together. The original
four-coordinate return-density analysis remains tied to the periodic
arithmetic and extracted critical edges developed below.

